\section{Discussion}\label{discussion}

In this work, we present a human MASLD Resource that records liver-trait genetic evidence and cross-sectional disease-state remodeling for each gene, with assay boundaries kept separate. Two earlier input defects—duplicate variant identifiers and non-PSD LD matrices across 90 primary eQTL fits, and repeated sequencing captures inflating the historical single-cell census—have been resolved across the synchronized figures and Unified Gene Catalog (yielding 430 primary multi-signal colocalized genes after allele-aware, PSD-projected refitting and 650,083 unique physical cells across 104 donors; the 117 frozen Hotspot program weights are unchanged). The bulk layer comes from 774 conservative genetic-identity clusters in five cohorts, not verified independent donors. After shared-posterior eligibility qualification and the PSD eQTL refit, the native multi-signal comparison includes 3,587 genes with a bulk test: 420 genetically supported, 209 canonical disease-state calls and 26 shared, against 20.9 expected within expression deciles (Mantel-Haenszel OR 1.321, 95\% CI 0.851 to 2.051; exact \emph{p} = 0.222). We detected neither enrichment nor depletion. Historical expression controls were not rerun in the qualified universe, so no direction is inferred (Supplementary Note 3). Another 935 disease-state calls have no multi-signal result, and 11 more are priority-untestable; these are not genetic negatives. The overlap establishes neither biological separation nor equivalence, independence or a privileged intersection. The separate signed comparison remains bounded: genome-wide rank agreement between ALT TWAS z and the disease contrast lies inside $\pm$0.03 without composition adjustment, but not with hepatocyte fraction in the design, and no direction is stated (Supplementary Note 5).

Most colocalized genes take their strongest evidence from liver-enzyme traits,
largely because far more gene-trait pairs were tested in that class. The
strongest native colocalization evidence comes from liver-enzyme traits for 417
of 430 genes and from direct disease or liver-fat traits for 13. Only 22 genes
have any direct-trait colocalization, and far more enzyme than direct-trait
gene-by-GWAS pairs were tested (12,121 against 771). Pooled across ancestries,
protein-altering plus canonical-splice variants carry 29.65\% of fine-mapped
posterior probability for direct MASLD and liver-fat traits and 6.67\% for enzyme
traits, counting credible sets at genome-wide-significant leads. This contrast is confounded with GWAS ancestry and fine-mapping
coverage, and no ancestry stratum supports a calibrated trait-class call
(Supplementary Note 4). A comparison of the trait classes within one ancestry
needs fine-mapping of all UKBB enzyme lead windows and of the MVP European ALT,
AST and NAFLD GWAS. Some protein-altering variants may act through protein
function without changing steady-state RNA. Others change RNA abundance through
nonsense-mediated decay or splicing (Rivas et al., 2015). One hypothesis for
which genes appear in both sets is therefore that expression-based
colocalization misses risk variants that act through protein function rather
than RNA abundance. The present comparison does not test it, because it does not
hold gene eligibility, phenotype, study power and expression observability
fixed. These results set one requirement for gene lists built this way in MASLD:
the trait that supplies each gene's evidence should be stated. A gene nominated
by ALT is not the same claim as a gene nominated by liver fat.

The historical named-r8 class comparison, not rerun after eligibility qualification, did not recover the prespecified prediction that
genetically supported genes are more hepatocyte-specific than disease-state
genes. In 281 expression-matched gene pairs, the observed difference was in the
opposite direction (Supplementary Note 3). Every test of that opposite direction
is post hoc, and none is calibrated. The difference comes mainly from the 56
pairs whose disease-state partner was itself tested by colocalization. We
therefore read it as a descriptive difference conditioned on selection, not as
confirmatory class separation or independent biological replication. The
protein-response comparison of 73 pairs gives no support for a class difference,
and it does not establish equivalence. Expression remains a major observability
difference between the classes. Among genes with a single-cell hepatocyte-specificity
value in the historical named-r8 classes (399 of 400 genetic-only and 1,124 of 1,129 state-only genes),
mean bulk expression of genetic-only genes exceeds that of state-only genes by
3.80 log2 units, and 281 of the 399 genetic-only genes obtain an
expression-matched partner. This 3.80 log2 gap belongs to the
hepatocyte-specificity comparison. Within the expression deciles of the
historical 3,769-gene overlap test, colocalized genes differ from other genes by 0.07 log2
units on average (0.70 before stratification). Unequal observability could affect which genes
appear in both sets, but the present comparison does not isolate that
explanation. In a prespecified test, colocalized genes were more
loss-of-function constrained (odds ratio 0.70 per SD of LOEUF, Holm \emph{p} 0.012),
consistent with the general pattern reported by Mostafavi et al.~(2023).
This model compares genetic-only with state-only genes and excludes the overlap;
it does not explain why genes occur in both sets. Germline regulatory effects may act before biopsy, in restricted
cell states or through small persistent shifts. Disease-state expression can
reflect injury, adaptation or composition. Neither class is automatically causal
or irrelevant.

The recurring examples make this distinction concrete. GNMT has a single-signal
direct-NAFLD colocalization (ABF PP.H4 0.706) but no multi-signal call. With
MAT1A and CYP2C19, it is lower in pooled RNA, in cross-sectional fibrosis
contrasts and in adjusted liver protein. In an exploratory analysis of 46 MASLD
participants, the three proteins covary. The covariation remains after
adjustment for fibrosis stage but not after adjustment for NAS or for all four
histology components (BH over 12 pair-by-adjustment tests). This pattern is
compatible with a shared association with histologic severity rather than a
three-gene circuit, but adjustment does not establish mediation. The ECM/IGFBP7 and ductular-injury/BICC1
programs show a different limit. The original atlas lacks ambient-RNA
correction, the ambient-RNA sensitivity analysis cannot resolve all datasets or
programs, and of the two only ECM/IGFBP7 exceeds a matched-gene spatial null,
which most of the 114 scorable fixed programs also exceed. Cross-assay
evidence therefore strengthens specific claims, but disease association, tissue
architecture, lineage origin and inherited mechanism remain distinct.

Each component of this Resource has precedent. Staged bulk liver transcriptomes
have been published (Hoang et al., 2019; Govaere et al., 2020; Pantano et al.,
2021), meta-analyzed with gene-level GWAS signal (Piras and DiStefano, 2024),
organized into network modules (Kumak et al., 2026) and ordered along a molecular continuum
(Kamzolas et al., 2026; \url{doi:10.1038/s42255-026-01543-7}). Single-cell and
spatial atlases resolve disease cell states (Ramachandran et al., 2019;
Guilliams et al., 2022; Boesch et al., 2026). A single-cell multiome assigns MASH
risk loci to liver cell types (Elison et al., 2025), single-cell eQTL mapping has been
done in MASLD liver biopsies (Hong et al., 2025), chromatin,
reporter-assay and CRISPR interference data assign target genes to MASLD risk
variants (Zhu et al., 2026), and large GWAS have combined
expression and coding evidence to nominate genes (Sveinbjornsson et al., 2022).
Multi-ancestry genetic association data have been used to nominate 212
candidate MASLD genes and, through genetically predicted expression, to select
drug-repurposing targets among them (Seagle et al., 2025).
This Resource adds a per-gene Catalog that keeps colocalization results beside
disease-state calls and records whether each layer was untested or tested
without support. Our reanalysis of a 138-donor liver caQTL map (Currin et al.,
2025) places a colocalized accessible region at 191 of 581 European liver-trait
GWAS signals that colocalized with a liver eQTL. At least 45 of these were
already reported by the source.

A related distinction between evidence roles has been reported at the protein layer. A
liver proteogenomic study of 504 biopsies found that differential protein
abundance and protein-QTL colocalization nominate largely different gene sets
(Gobeil et al., 2026b). Of the
five genes that study nominated, our map places GCKR on the primary
expression-colocalization side, MTTP colocalizes only in the
single-signal (ABF) sensitivity analysis (best PP.H4 0.82), ERLIN1 does not
reach that set (best ABF PP.H4 0.47; Catalog class unresolved), and LGALS1 is in the established-state class. MTARC1 lies outside our
expression map but is resolved by the protein-QTL analysis of that study. GCKR
has differing expression-QTL and protein-QTL directions. These cases show the
difference between protein and expression observability. They do not establish
a binary split into causal and reactive genes. A single-cell multiomic liver
study likewise reports that chromatin-accessibility QTL and expression QTL
colocalize with largely distinct risk loci. It supplies cell-type regulatory
context that tissue-level eQTL analysis cannot provide (Elison et al., 2025).
These external studies extend the assay boundary that the Gene Catalog records.
They are credited source comparisons, not results of this paper.

The gene and module results of Kumak et al.~are likewise source findings, not
results of this paper. Of their NAS-associated and fibrosis-associated genes,
944 and 1,013 map to the Catalog, and 258 and 301 of these reach the canonical
bulk call here. Of those present in both the bulk and colocalization tables, 661
of 902 and 714 of 943 had no multi-signal colocalization result. The Catalog
keeps this state separate from a tested negative. Among the 241 NAS and 228
fibrosis genes that were evaluated, 36 and 28 colocalized, against 28.9 and 24.9
expected for genes of the same expression level (Mantel-Haenszel OR 1.32, 95\% CI
0.91 to 1.92, and OR 1.16, 95\% CI 0.76 to 1.76; Holm-adjusted \emph{p} 0.29 and 0.51,
Holm over the two lists). Thus, among genes evaluated by multi-signal COLOC, the
colocalization rate of these genes did not differ detectably from that of genes
of the same expression level (power in Supplementary Note 15). Kumak et al.
prioritize genes by combining meta-analysis evidence with centrality in a
curated protein interaction network. The 117 programs and 43 systems used here
were fixed before any outcome was read and are never reweighted by network
position. A sparsely annotated gene therefore keeps its assigned weight, and an
inadequately observed gene is recorded as untestable rather than ranked low.

The 1,854 pooled disease-state calls describe tissue associations in the
discovery cohorts, not an independently replicated disease signature. The
current random-effects sensitivity retains the pooled direction for 1,850 genes
and one direction in every available cohort for 1,208, but none passes BH over
24,115 genes with modified Hartung--Knapp uncertainty. This combination is
consistent with broad directional agreement and substantial uncertainty about
effect magnitudes across four/five cohorts. It does not establish absence or
equivalence. The earlier ``consistent'' classifications used different inputs
and uncertainty rules and do not certify replication. SteatoSITE provides a
separate test of cross-sectional fibrosis-stage effects; its success cannot be
substituted for disease-versus-control replication. The measured tissue
associations remain useful for follow-up, with their cohort and contrast limits
attached to each interpretation.

The stage-stratified analyses add biological context without establishing
change over time. Adjacent-stage contrasts and cell-resolved programs describe
stage-ordered remodeling across different participants. The externally anchored
molecular continuum, adapted from Kamzolas et al.~(2026), adds a separate
molecular view without redefining disease. It aligns strongly with the
fibrosis-based pathway, NMF and Hotspot results, which is consistent with the
major biological programs being preserved. This alignment was computed on the
same data and is not independent validation. Stage-adjusted associations show
that molecular variation within histologic bins tracks the continuum, and which
programs carry that variation depends on adjustment for cell-type composition.
The prespecified primary continuum test missed in both test cohorts (Holm \emph{p}
0.082 in GSE162694 and GSE213621), and matched random gene sets give a mean of
86 continuum-associated programs against the 92 reported (\emph{p} 0.109).
Histology therefore remains the interpretable clinical anchor, and the continuum
adds molecular resolution within a stage. All 117 Hotspot program definitions
and 43 system memberships remain fixed. An outcome-blind map extends this
comparison across complete pathway and Hotspot layers. Stage and continuum
effects align broadly, but only a minority of systems keep one sliding-window
direction across all five cohorts. This supports shared but cohort-variable
remodeling, not a universal molecular sequence or ontology. Program-pair
covariance is reported in Results. It could not be tested in GSE213621, where
the program-group permutation control failed. Bulk covariance also mixes cell
abundance with program activity, so it does not show coordination between cell
types. Neither analysis identifies transitions within individuals, a dominant
fibrosis switch or therapeutic windows. Those claims need repeated sampling of
the same participants and perturbation experiments.

For users of the Resource, the main lesson of Figure S7 is that each model score
must be read beside its exposure audit and its task-native baseline.
Our per-cohort exposure audit showed that all 109 GSE268273
validation samples were in the BulkFormer pretraining data, so that model could not
be evaluated on samples it had not seen. Against task-native baselines, the paired
difference between the severity instrument and the in-cohort ridge has a 95\%
interval that includes zero, and GSE268273 was also used to choose the
expression transform. A per-region model of H3K27ac that adds the expression of genes near each
region (the cis head) scored higher than a model on global expression components
alone, but its gain over a reduced-rank RNA model like that of Figure S7B was
smaller (Supplementary Note 14).

There are several limitations to this study. The liver-eQTL panel is
reused and of European ancestry, so non-European GWAS lack ancestry-matched
molecular QTL data. The regulatory-genetic layer is based on expression QTL, so
protein-altering and protein-level mechanisms that leave the transcript
unchanged, such as MTARC1, fall outside its resolution. This layer is also
tissue-level, not cell-type-resolved. Source-reported protein-QTL and
single-cell chromatin studies supply distinct assay contexts, but no single
layer is a complete causal filter. Most genetically supported genes come from
liver-enzyme rather than direct MASLD traits. Partitioned polygenic scores
already show that MASLD risk variants do not form one class (Jamialahmadi et
al., 2024). The direct-versus-enzyme split used here is a different, trait-based
partition. The class supported on both axes is small, and every context assay
has source-specific coverage. No accepted spatial analysis provides a
MASLD-versus-healthy disease-effect estimand, and this paper contains no new
experimental perturbation. The GTEx comparison measures histologic fat of
unrecorded cause in postmortem liver. Its slide tissue-area negative control
missed its prediction (Spearman \emph{rho} 0.354, \emph{p} = 0.0564) and came close to its
stop rule. In a post hoc check, adding tissue area to the model left the fat
correlation at 0.631 (0.638 without it). Ischemic time was also correlated with the same genes' disease
effects, with the opposite sign and less strongly than fat (\emph{rho} −0.350 against
0.636). The caQTL reanalysis rests on one 138-donor map. Its chr1 and chr19
positive controls were seen in a pilot and are not blind, and at least 45 of the
191 colocalized signals were already reported by its authors. We did not check
overlap with published three-way eQTL, caQTL and GWAS colocalizations
(Broadaway et al., 2024). We record these as Catalog limitations and do not
infer mechanisms from source-owned data.

Further limitations concern units of analysis and models. Genotypes called from RNA-seq
reads identified repeated libraries and conflicting labels, but these calls,
sample-key checks and expression-based QC do not give complete genotype-based
identity confirmation, which was not established here. The bulk units therefore remain genetic-identity
clusters, not verified donors. Gene-level totals can hide opposing isoform
changes. Differential gene expression concerns total gene abundance,
differential transcript expression concerns individual transcript abundance, and
differential transcript usage concerns relative isoform proportions within a
gene. This Resource tests the first, not isoform switching. Short reads shared
among isoforms add assignment uncertainty, and transcript estimates depend on
the annotation. Oakes and Pachter (2026) emphasized these distinctions and the need to
separate demonstrated responses from exploratory physiological explanations. Model validation is specific to its task, cohort and
training exposure, and rank correlation is not absolute-stage calibration. The
GSE193066 repeat-biopsy intraclass correlation (ICC) and SD 0.6433 use
source-declared pairs whose identity was not qualified. They cannot establish
primary within-person reliability or label-error calibration, and ceiling values
derived from them remain sensitivity scenarios. Serial biopsies can include
shared technical effects and true change, and uncertainty conditional on fitted
predictions omits retraining variation. The comparator models were fitted
within each test cohort, so they had already seen that cohort's data.
Comparison with them therefore cannot show that the models perform equally on
unseen cohorts, or that one model class is better in general.

Two questions stay open. We do not assign genes or programs separately to steatosis, inflammation or fibrosis; fibrosis stage is the severity anchor, and NAS enters only as a sum or as an adjustment set. Colocalization here tests single variants, so the Resource does not test whether combinations of variants act differently, and the local haplotype at GNMT is unresolved. Both questions are recorded as indeterminate, not negative.

The chromatin analyses further distinguish disease accessibility from predicted
regulation. Modest cross-cohort agreement and four stellate peakwise candidates
do not establish a replicated MASH regulatory program, particularly when
recorded-batch adjustment substantially changes discovery counts. Cicero's
small held-donor hepatocyte RNA benefit survives the fixed unshuffled stress
tests, but its advantage over proximity also survives RNA-pairing shuffles.
These data support descriptive regulatory context without validating individual
edges or a pairing-specific mechanism. CHST9 and TMEM150B are single-cohort
hepatocyte RNA candidates in both recorded-batch and condition-only models;
their expression associations do not establish that a linked accessible region
regulates either gene. Missing qualified cross-assay donor matches and
insufficient normal donors in the other RNA lineages remain limits on what can
be tested (Supplementary Note 16).

Within these limits, the Resource gives each gene an experiment matched to its evidence. Genetically anchored candidates motivate allele-to-gene and
context-appropriate perturbation without requiring disease-state expression.
Disease-state-associated candidates motivate tests of active tissue function and
reversibility without being labeled inherited drivers. Candidates supported on
both axes motivate experiments that connect the two relationships in one model.
Untested and unresolved candidates, 26,688 of the 28,233 genes (13,733 untested and 12,955 unresolved), first need the
missing decisive measurement. The Gene Catalog and the source tables needed to check each call
are deposited with this paper (Data availability), so that others can look up
these distinctions for any gene instead of relying on a single aggregate rank.
